%==============================================================================
\section{Supervised vs Unsupervised Learning --- So sánh hai hướng học}
\label{sec:supervised-vs-unsupervised}
%==============================================================================

\begin{ynghia}[title={Nguồn}]
\tsurl{https://www.tutorialspoint.com/machine_learning/machine_learning_supervised_vs_unsupervised_learning.htm}.
\end{ynghia}

\begin{dongco}[title={Hai hướng phổ biến nhất}]
Supervised và unsupervised learning là hai cách tiếp cận phổ biến trong ML.
Cách phân biệt đơn giản nhất: \textbf{loại tập huấn luyện} và \textbf{cách huấn luyện mô hình} --- nhưng còn nhiều khác biệt khác, trình bày dưới đây.
\end{dongco}

\begin{luuy}[title={Đối chiếu}]
Hai bài \textbf{Supervised Learning} và \textbf{Unsupervised Learning} đi sâu từng hướng; bài này tập trung \textbf{bảng so sánh} và \textbf{chọn phương pháp}.
\end{luuy}

\subsection{Supervised Learning là gì? (nhắc lại)}

\begin{dinhnghia}[title={Học có giám sát}]
Dùng dataset \textbf{có nhãn} để huấn luyện mô hình --- phù hợp phân loại hoặc dự đoán đầu ra.
Hai nhóm tác vụ:
\end{dinhnghia}

\textbf{1. Classification} --- dự đoán giá trị phân loại (spam/không spam, đúng/sai, \ldots).
Thuật toán học ánh xạ đầu vào $\rightarrow$ nhãn lớp.
Ví dụ: KNN, Random Forest, Decision Trees.

\textbf{2. Regression} --- phân tích quan hệ giữa điểm dữ liệu, dự báo giá trị số liên tục (giá nhà, doanh số, \ldots).
Ví dụ: linear regression, polynomial regression, logistic regression (nguồn liệt kê logistic trong regression --- thường dùng cho classification; giữ theo nguồn).

\subsection{Unsupervised Learning là gì? (nhắc lại)}

\begin{dinhnghia}[title={Học không giám sát}]
Huấn luyện trên dữ liệu \textbf{thô, không nhãn} để tìm pattern không cần giám sát con người.
\end{dinhnghia}

\textbf{1. Clustering} --- gom điểm dữ liệu theo độ tương đồng (ví dụ K-means).

\textbf{2. Association} --- nhóm theo quy tắc định sẵn; dùng trong Market Basket Analysis (Apriori).

\textbf{3. Dimensionality Reduction} --- giảm kích thước dataset bằng cách bỏ feature không cần thiết mà vẫn giữ bản chất dữ liệu.

\subsection{Bảng khác biệt chính}

\begin{vidu}[title={Supervised vs Unsupervised (dịch từ bảng nguồn)}]
\begin{tabularx}{\linewidth}{>{\raggedright\arraybackslash}p{2.8cm}XX}
\textbf{Tiêu chí} & \textbf{Supervised Learning} & \textbf{Unsupervised Learning} \\
\midrule
Định nghĩa &
Thuật toán huấn luyện trên dữ liệu mà mỗi đầu vào có đầu ra tương ứng. &
Thuật toán tìm pattern trong dữ liệu không có nhãn định sẵn. \\
Mục tiêu &
Dự đoán hoặc phân loại dựa trên feature đầu vào. &
Khám phá pattern, cấu trúc và quan hệ ẩn. \\
Dữ liệu đầu vào &
Có nhãn: đầu vào kèm nhãn đầu ra. &
Không nhãn: dữ liệu thô, chưa gán nhãn. \\
Giám sát con người &
Cần giám sát để huấn luyện (gán nhãn). &
Không cần giám sát khi huấn luyện. \\
Tác vụ &
Regression, Classification. &
Clustering, Association, Dimensionality Reduction. \\
Độ phức tạp tính toán &
Thường đơn giản hơn (theo nguồn). &
Thường phức tạp hơn (theo nguồn). \\
Thuật toán &
Linear regression, KNN, Decision Trees, Naive Bayes, SVM. &
K-Means, DBSCAN, Autoencoders. \\
Độ chính xác &
Thường chính xác cao hơn khi có nhãn đúng. &
Thường kém chính xác hơn (theo nguồn). \\
Ứng dụng &
Phân loại ảnh, sentiment analysis, hệ gợi ý. &
Phân khúc khách hàng, anomaly detection, recommendation engines, NLP. \\
\end{tabularx}
\end{vidu}

\begin{luuy}[title={Giảng viên lưu ý}]
Cột ``Độ phức tạp tính toán'' và ``Độ chính xác'' trong bảng nguồn là \textbf{khái quát hoá} --- phụ thuộc thuật toán, dữ liệu và metric. Unsupervised khó so ``chính xác'' vì thường không có nhãn chuẩn. Logistic regression thuộc \textbf{phân loại}, không phải hồi quy số.
\end{luuy}

\subsection{Chọn Supervised hay Unsupervised?}

\begin{phuongphap}[title={Các yếu tố cần xem xét}]
\begin{itemize}
  \item \textbf{Dataset}: có nhãn hay không? có đủ thời gian, nguồn lực, chuyên môn để gán nhãn?
  \item \textbf{Mục tiêu}: phân loại, khám phá pattern mới, hay mô hình dự báo?
  \item \textbf{Thuật toán}: số feature, số chiều, khối lượng dữ liệu có phù hợp không?
\end{itemize}
\end{phuongphap}

\subsection{Semi-supervised Learning (lối thoát trung gian)}

\begin{ynghia}[title={Khi phân vân giữa hai hướng}]
Semi-supervised là ``điểm giữa an toàn'': kết hợp supervised và unsupervised --- \textbf{phần nhỏ} có nhãn, \textbf{phần lớn} không nhãn.
Phù hợp khi có rất nhiều dữ liệu nhưng khó xác định feature liên quan thủ công.
Chi tiết ở bài \textbf{Semi-Supervised Learning}.
\end{ynghia}
